\chapter{Research Development and Design Decisions}
\label{ch:development}

\section{From component design to complementary evidence}
The early graph and temporal experiments showed that relationships among observations helped detection. Inspection of the expert outputs then exposed weaknesses in the reference and routing. The later work addressed these problems through residual modelling, temporal specialists and checks on the alarms they added to the system.

The experiments served different purposes. Historical controls compared components within their own benchmark; operational development and calibration guided the design and operating points. Frozen confirmation then evaluated the selected designs. Appendix~\ref{app:evidence} reproduces the quantitative evidence and identifies the accompanying experiment records for each level.

Figure~\ref{fig:development} summarizes the questions that guided the architectural development. Rejected alternatives are discussed in the following sections. Drift, noise and clean-period false alarms drove the main architectural changes; replay later tested the adequacy of General's representation in L-Town.

\figfile{development.pdf}{1.0}{Questions guiding architectural development}{Conceptual development of the research. Each stage identifies a question that changed the design. Arrows indicate the reasoning sequence, not a controlled increase in performance across identical datasets. Both Modena and L-Town belong to the final evaluation.}{fig:development}

\section{Starting with graph and temporal models}
The initial neural design encoded sensor information on the network graph and processed it over time. A shared spatial backbone represented each snapshot; recurrence combined representations across a window. Reconstruction outputs and anomaly heads then used the resulting features. The architecture was intended to use neighbouring observations and temporal context to interpret a pressure reading.

Matched controls separated spatial and temporal contributions. The topology-free model disabled graph messages while preserving the temporal path. The one-step temporal ablation used endpoints and validation conditions aligned with those of the longer-window model.

The historical architecture in Figure~\ref{fig:architectures} extended this design by assigning components to clean operation and the five attack families. Each expert retained a full spatio-temporal predictor, and the router combined outputs. The six-expert mixture therefore added six full predictors, not just small output heads. The comparison with a single temporal model tested whether this extra capacity supplied useful information and whether the router retained it.

\figfile{architectures.pdf}{1.0}{Historical and selected architectural organization}{Two architectural stages at different levels of abstraction. The historical system mixes six graph-temporal experts. The selected hybrid has three top-level branches, with five internal components inside General. The diagram describes organization, not an isolated same-data performance comparison.}{fig:architectures}

The historical episode ablation preserved ten fits for each of four configurations on one generated Modena dataset. The stored results showed that removing topology or temporal context had a larger F1 cost than removing the mixture. Because all forty fits shared one underlying dataset, the next confirmation varied data seeds as well as model seeds [\evid{E4}].

\section{What the controlled episode experiments established}
The spatial--temporal confirmation compared matched spatial and temporal models over three data seeds, with three model seeds nested within each. The mean AUPRC gain of the temporal model was positive in both networks. The topology comparison likewise retained temporal modelling while disabling messages in the control. Its positive differences supported the use of relational information in those episode experiments [\evid{E4}].

\begin{table}[tbp]
\centering\small
\caption[Historical controlled episode comparisons]{Historical controlled episode comparisons. Entries are mean paired AUPRC differences, not final hybrid scores. Reported 95\% intervals use three data-seed means; each data seed contains three paired model fits. Source: preserved episode confirmation summaries, reproduced in Appendix~\ref{app:evidence}.}
\label{tab:historical}
\begin{tabular}{@{}llrr@{}}
\toprule
Comparison & Network & Mean $\Delta$AUPRC & 95\% interval\\
\midrule
Temporal minus spatial & Modena & +0.0798 & [0.0707, 0.0889]\\
Temporal minus spatial & L-Town & +0.0890 & [0.0514, 0.1267]\\
Topology-aware minus free & Modena & +0.0850 & [0.0395, 0.1305]\\
Topology-aware minus free & L-Town & +0.0820 & [0.0432, 0.1208]\\
MoE minus temporal & Modena & +0.0335 & [0.0043, 0.0626]\\
\bottomrule
\end{tabular}
\par\smallskip\begin{minipage}{\linewidth}\footnotesize\textit{Scope note.} The retained matched MoE--temporal confirmation covers Modena only (three data seeds, three model seeds each). A corresponding L-Town confirmation is absent from this campaign; no L-Town MoE effect is estimated here.\end{minipage}
\end{table}

Topology and temporal context contributed in both networks under the controlled episode protocol. The mixture added a smaller positive average gain in Modena. The intervals summarize three data-seed means; the table's scope note identifies the absent L-Town mixture comparison. Figure~\ref{fig:effects} displays the individual data-seed means.

\figfile{historical_effects.pdf}{0.96}{Historical paired AUPRC differences by data seed}{Data-seed mean AUPRC differences in the historical controls. Each point averages paired model fits for one data seed; the horizontal axis is a paired change. The comparisons share the same scale. The figure covers the matched historical controls.}{fig:effects}

The matched temporal-permutation control examined the role of coherent episode structure. In Modena, the mean temporal AUPRC gain was approximately 0.0757 on episodes and 0.0401 in the matched IID control, with a mean interaction of 0.0357. Persistence thus contributed to the measured advantage, while the positive IID gain showed that recurrence also exploited other temporal information.

This distinction motivated explicit representations of temporal change in the later design: the useful signal and its observation window could then be inspected directly.

\section{Expert labels did not settle expert quality}
The neural mixture required its own diagnostic beyond the aggregate MoE comparison. The Modena specialization analysis compared soft mixing with hard top-1 routing and a diagnostic using true family information. The mean soft-mixture AUPRC was 0.5637, while mean router accuracy was 0.4014. The family-owning expert ranked first in 40 of 45 family-by-seed comparisons, but specialization was incomplete and hard top-1 performance was weaker than soft mixing [\evid{E4}].

The diagnostics suggest useful expert information that the router does not reliably isolate. True-family routing helps investigate this mismatch, but requires labels unavailable at inference; it is neither a deployable candidate nor a universal upper bound, since even the family-owning expert can make endpoint errors.

Correctly recognizing a family does not ensure that the router selects the best detector for a reading. An expert may also discriminate its own family well while producing false positives elsewhere. Reporting only its assigned-family performance would miss these problems.

Specialization was therefore assessed by the true detections it added beyond General and the errors it introduced elsewhere. This criterion guided the subsequent Drift and Noise branches.

\section{Changing the normal reference}
In the earlier reconstruction path, a reading could enter the representation used to predict its own expected value. This allowed corruption to affect the reference against which it was judged.

The blind-reference prototype prevented this direct copying by predicting each target group from observed sensors outside it. It fitted a low-rank representation to noisy observed pressures in normal training episodes. The target group also remained excluded during robust reweighting [\evid{E5}].

The prototype supplied tree classifiers with normalized residual magnitude and temporal summaries such as slope, variability and persistence, together with observation support. It also introduced five internal components: general, abrupt, replay, drift and noise. Random and targeted corruption shared the abrupt mechanism in this internal decomposition, while their evaluation labels remained separate. This organization is the predecessor of the final General branch.

In the small development validation comparison, the blind-reference mixture performed strongly on random and targeted attacks, but drift and noise F1 were only about 0.2222 and 0.2985. The aggregate improvement justified further work on the reference, while these weak families required a different approach.

The prototype changed both the representation and estimator, so their separate contributions cannot be inferred from this gain. The earlier matched neural controls provide the evidence for graph and temporal contributions.

A later control replaced the reference-derived inputs while holding the fitted dynamic-feature classifier and its thresholds fixed. Pooled F1 changed from 0.8181 to 0.8335 at the overall-priority threshold, but from 0.7325 to 0.7115 at the family-balanced threshold. This comparison evaluates the multi-view, noise-scaled and leverage-adjusted reference as a package. Both references already excluded the target group and used robust reweighting, so the comparison does not isolate the effects of exclusion, robustness or their interaction. The thresholds were retained from the earlier input distribution; the comparison does not estimate the performance of separately recalibrated systems [\evid{E5}]. Section~\ref{sec:reference_factorial} defines a separate factorial comparison of target-group exclusion and iterative reweighting, with downstream refitting and a common calibration policy. It holds whitening and median aggregation fixed, so its scope is narrower than robustness as a whole.

\section{Developing Drift and Noise evidence}
Gradual bias and excess noise posed different temporal problems. For drift, the early part of a ramp can remain close to the reference, so evidence may need to accumulate across time. For noise, individual values can still look ordinary even when their local distribution has changed. General residual magnitude alone did not give the desired coverage, and similarly structured family experts were not reliably better than the general component.

The next experiments represented directional consistency, dispersion, reference support and daily context explicitly. They led to the seasonal and fast-response components. Daily comparisons reflect the generator's demand pattern, so their usefulness depends on that temporal context.

Other robust-reference experiments changed classifier training, weighting or score combination at the same time. Their results concern those combined changes, unlike the fixed-classifier reference control above.

Alarm persistence caused a further problem. Retaining a positive state could help detect weak episodes but also prolong alarms after an event ended, even when family F1 looked favourable. Candidate temporal mechanisms were therefore checked for clean-period false positives and their ability to return to normal after an event.

Drift and Noise became separate branches because they required temporal evidence and treatment distinct from the broad residual path. General retained its role in detecting abrupt, targeted and other observable discrepancies.

\section{Making latency part of the task}
The delay study allowed three hours for specialist evidence to accumulate. It recorded both the endpoint time and the latest observation time permitted for that decision.

An early screening study used the maximum of existing specialist scores over a bounded future window. Its purpose was to determine whether later evidence could improve the weak-family ranking under a fixed latency allowance. In that TRAIN-based hindsight screen, the forward maximum improved the reported drift and noise F1 relative to the zero-delay score. A learned delayed head sometimes improved detection scores but also kept producing false positives after an event ended, beyond the accepted limit [\evid{E6}].

Following the TRAIN hindsight screen, a Modena study compared separately calibrated zero-hour and three-hour systems on six held evaluation sources, comprising 144 scenarios and 3,015,103 observed pressure endpoints. Both operating points were frozen on calibration before evaluation, under the same clean-FPR ceiling of 0.005. Table~\ref{tab:historicaldelay} reproduces that comparison from the retained records [\evid{E6}].

\begin{table}[!htbp]
\centering\small
\caption[Historical Modena comparison of zero and three hours]{Historical Modena latency comparison: one fitted model configuration, six shared evaluation sources. F1 is computed over the concatenated endpoints, rather than averaged across the final campaign's model/source cells. Clean FPR is shown as a percentage; its difference is in percentage points.}
\label{tab:historicaldelay}
\begin{tabular}{@{}lrrr@{}}
\toprule
Metric & Zero hours & Three hours & Difference\\
\midrule
Random F1 & 0.9359 & 0.9427 & +0.0068\\
Replay F1 & 0.7955 & 0.7886 & $-0.0069$\\
Drift F1 & 0.7618 & 0.8368 & +0.0750\\
Noise F1 & 0.7346 & 0.8437 & +0.1091\\
Targeted F1 & 0.9542 & 0.9605 & +0.0064\\
Pooled F1 & 0.6812 & 0.6867 & +0.0054\\
Clean FPR (\%) & 0.3126 & 0.3847 & +0.0721\\
\bottomrule
\end{tabular}
\end{table}

The delayed configuration improves drift and noise substantially, while replay declines slightly and clean false alarms increase. General remains immediate in both configurations. At three hours, the specialists blend a forward maximum with a delayed head; the zero-hour control uses causal scores. Calibration also selects different integration rules. This comparison therefore measures the benefit and cost of permitting a delayed configuration, including its recalibration, rather than the isolated effect of waiting three hours. It supports the development decision to use later evidence for weak temporal families, but does not establish that three hours is optimal or provide a latency curve for the final ten-seed systems.

The retained design therefore keeps General immediate and finalizes specialist decisions on the three-hour clock. Chapter~\ref{ch:architecture} specifies how the bounded observations enter that decision.

\section{From expert outputs to guarded decisions}
Accepting every specialist alarm also admitted errors caused by other mechanisms. Evidence routing and local feedback were developed to control those errors across the evaluated corpus.

Routing describes the mechanism suggested by network-level evidence. Local feedback and verification then assess whether a specialist alarm should be retained. Their contribution is measured through final detection errors, with General's decision path preserved.

An earlier evaluation of the guarded system is recorded as EVAL-3 in the project archive. Its matched comparisons established the value of regulating specialist decisions in that evaluation population. The subsequent ten-model-seed campaign assessed the later selected implementation on its own declared sources; Appendix~\ref{app:evidence} keeps the two records separate.

The subsequent matched integration campaign (Stage-E in the project records) compared the baseline with both a threshold-only alternative and a candidate adding verification and early-warning evidence. This separates the benefit available from recalibration alone from that of the larger intervention [\evid{E1}].

Section~\ref{sec:operating_sensitivity-supp20261008} examines the conditional contribution of each branch and guard to the final alarm rule. Those supplementary removals retain the fitted models and original thresholds, so their effects differ from those of retraining or recalibrating a reduced architecture.

The decision heads also required out-of-fold training. A verifier trained only on in-sample expert outputs could learn from unrealistically clean score patterns. Source-held expert predictions provide a more relevant training signal for the integration layer. In the final extension, the stochastic components were refitted for each added model seed rather than varying only the last threshold or router parameter. Table~\ref{tab:decisions} summarizes the design decisions and their evidence boundaries.

\begin{table}[tbp]
\centering\small
\caption[Design decisions and the observations that motivated them]{Design decisions and the observations that motivated them. The evidence column distinguishes controlled comparisons from development studies; the latter explain the choices without isolating causal effects.}
\label{tab:decisions}
\begin{tabularx}{\linewidth}{@{}YYY@{}}
\toprule
Observation & Architectural response & Evidence boundary\\
\midrule
Topology and recurrence contributed under episode controls & Retain spatial/temporal lessons in the design rationale & Historical matched comparisons\\
A named expert was not always the best detector & Inspect component and mixture outputs separately & Specialization diagnostics\\
A target observation could influence its own reconstruction & Use target-group-excluding reference predictions & Reference prototypes and controls\\
Drift/noise remained weak under broad residual evidence & Develop specialized temporal representations & Operational development\\
Unconditional specialization could add false positives & Retain General and regulate specialist alarms & Matched integration comparisons\\
Later evidence helped weak-family discrimination & Declare specialist decision latency & TRAIN screen followed by calibrated evaluation\\
\bottomrule
\end{tabularx}
\end{table}

\section{A bounded refinement for a difficult family}
\label{sec:flowcandidates}
L-Town replay tested whether the established architecture needed a different operating point or a more informative representation. Acceptance still required balanced performance across the complete system.

A symmetric recalibration of the internal mixture router improved replay in its fresh paired confirmation, from mean F1 0.6461 to 0.6757, while pooled F1 increased from 0.8799 to 0.8841. A preliminary attempt to force a larger false-alarm budget had produced little replay benefit and a large pooled loss. These outcomes justified examining representation rather than assuming that a more permissive threshold would solve the remaining problem [\evid{E7}].

The first shared-history candidate added 56 features summarizing histories of existing residual evidence. On the held-source TRAIN diagnostic, General-mixture replay F1 fell from 0.7173 to 0.6954, leading to rejection.

The subsequent representation used differences between actually received pressures at common lags, with explicit availability indicators. Fourteen lags supplied signed change, absolute change and availability, giving 42 new features. The same bank was made available to the five internal General components and their router; it did not add a new top-level expert or a replay-specific inference rule. Differences were placed in 0.1-reference-scale bins to reduce reliance on exact floating-point equality [\evid{E8}].

The initial TRAIN pilot improved General-mixture replay F1 to 0.7432, but it contained only two replay events. The broader TRAIN replication improved replay in all 20 primary source-held fold/seed comparisons while failing the declared noise-protection gate. Its mean noise change was slightly below the permitted limit, and one source showed a larger loss. The candidate was therefore rejected at this stage.

A separate protected full-system experiment subsequently assessed the representation with the external specialized branches retained. It compared predeclared General-update variants on calibration, selected full history and froze the design before fresh confirmation. The five-model-seed confirmation improved replay from its paired baseline of 0.7064 to 0.7303, with protection checks passing. Appendix~\ref{app:evidence} supplies all 30 paired replay values and their calculation. This comparison uses its own paired baseline, distinct from the earlier 0.6757 measured on other confirmation sources [\evid{E8}].

\noindent\textbf{Flow-based development candidates.}
Flow information was tested in two development stages. An earlier Modena study used source 811 to compare local pressure features, pressure-context features and joint pressure/flow context for the weak-family experts. These three feature candidates used matching classifier hyperparameters, with references fitted inside their training folds. Adding flow did not improve normal-reference accuracy: mean absolute error increased from 0.1892 to 0.2043 m on calibration and from 0.1678 to 0.1810 m on the reused development validation set. The joint candidate reached pooled validation F1 of 0.8484, but drift and noise F1 were 0.2593 and 0.3429. Calibration selected the pressure-context candidate; the reused validation set did not trigger reselection. That set contained only one event for each weak family, limiting the finding to this historical development comparison [\evid{E10}].

The later L-Town replay study compared the received-pressure-history reference with three alternatives: internal probability calibration and routing correction, received-flow features, and their combination. The flow representation added 44 inputs to the existing 158: a rank-16 masked flow representation, a regularized regression from that representation to normal pressure, pressure discrepancies, availability measures, and pressure/flow change consistency at 1, 3, 6, 12 and 24 hours. All references and teachers excluded the relevant held sources. Five source folds and five model seeds produced 25 paired General-branch diagnostics.

Mean TRAIN replay F1 was 0.704881 for the reference, 0.703202 for probability fusion, 0.718817 for flow and 0.706480 for their combination. The three candidates improved on the reference in 12, 23 and 14 of the 25 comparisons, respectively. Eligibility required at least 20 wins, mean replay gain of at least 0.01 and a positive mean gain in every source. Only the flow candidate qualified, and it also passed the limited pressure-feature rounding checks. At full-system calibration, however, none of the five fits had an operating point satisfying the protection constraints. At the rejected maximum-worst-family settings, all five seeds failed the strict requirement that replay F1 exceed its paired reference (deficits of 0.002410--0.003376). Seeds 702 and 705 also violated the no-loss pooled-F1 and no-increase clean-FPR constraints; seed 702 additionally exceeded the 0.005 permitted F1 loss in each of the other four attack families, whereas seeds 701, 703 and 704 passed the remaining aggregate protections. The rejected settings that maximized worst-family F1 averaged replay F1 of 0.702639, compared with 0.705392 for their paired reference. The candidate did not proceed to fresh confirmation [\evid{E10}].

A subsequent L-Town candidate combined the same 44 flow inputs with 140 received-trajectory additions, giving 342 inputs in total. Its General-branch TRAIN replay F1 increased from 0.704881 to 0.738154, with improvement in all 25 comparisons. Full-system calibration satisfied the family, pooled-performance and false-alarm protections, but replay increased only from 0.705392 to 0.709936. The paired gain of 0.004544 fell short of the predeclared 0.02 progression requirement, preventing fresh confirmation. These calibration values use a different corpus from the confirmed 0.730307 history result and cannot be added to it. The confirmed pressure-history system remained the selected L-Town design [\evid{E10}].

\section{Settling the two-network architecture}
Modena retains the confirmed Stage-E system; L-Town retains its received-pressure-history version of General. Their shared three-branch framework combines broad evidence, temporal specialization and controlled specialist acceptance. Chapter~\ref{ch:architecture} specifies both implementations, and Chapter~\ref{ch:results} evaluates their performance.

The final extension increased model seeds from 701--705 to 701--710 on each network. It preserved the attack distribution and selection procedure, freezing the added fits and calibration rules before evaluation on the existing six sources [\evid{E2}].
